% Source: https://artofproblemsolving.com/community/c3h1922949_hole_vs_asymptote
% Stable identity: aops-topic-1922949-q01
\textbf{(a)} The denominator factors as $x^2-3x+2=(x-1)(x-2)$, so the domain is all real $x$ except $\boxed{x=1\text{ and }x=2}$.

\textbf{(b)} The numerator is $(x-1)(x+1)$, so for $x\ne1,2$,
\[
f(x)=\dfrac{(x-1)(x+1)}{(x-1)(x-2)}=\dfrac{x+1}{x-2}.
\]
\emph{At $x=1$:} $f(1)$ is undefined, but near $x=1$ the values $f(x)=\dfrac{x+1}{x-2}$ approach $\dfrac{2}{-1}=-2$, a finite number. The graph is the curve $y=\dfrac{x+1}{x-2}$ with the single point $(1,-2)$ missing: a \textbf{hole} at $\boxed{(1,-2)}$.

\emph{At $x=2$:} the numerator $x+1$ approaches $3\ne0$ while the denominator $x-2$ approaches $0$, so $|f(x)|$ grows without bound ($f(x)\to+\infty$ as $x\to2^+$ and $f(x)\to-\infty$ as $x\to2^-$). The line $\boxed{x=2}$ is a \textbf{vertical asymptote}.

\textbf{(c)} Factor the numerator and denominator and cancel common factors as far as possible. If $x=c$ makes the original denominator $0$ but the factor $(x-c)$ cancels completely from the denominator, the graph has a hole at $x=c$ (its $y$-coordinate is the value of the reduced expression at $c$). If $(x-c)$ is still a factor of the reduced denominator, the graph has a vertical asymptote $x=c$.
